\section{Introduction}
\label{sec:introduction}

Automatic summarisation is evaluated far more often than humans can rate it. Every training run, model comparison and ablation needs a score for thousands of candidate summaries, so the field relies on automatic metrics, and those metrics are trusted because they correlate with human judgement on a few meta-evaluation benchmarks. SummEval~\citep{fabbri2021summeval} is the most widely used of them: 100 CNN/DailyMail articles~\citep{hermann2015cnndm}, 16 summarisation systems, and expert ratings of coherence, consistency, fluency and relevance for every summary.

The metrics that agree best with these ratings are expensive or need resources that are often missing. Reference-based metrics, from ROUGE~\citep{lin2004rouge} to BERTScore~\citep{zhang2020bertscore}, need a human-written reference summary. Learned evaluators such as UniEval~\citep{zhong2022unieval} need a reference for some dimensions, and large language model (LLM) judges such as G-Eval~\citep{liu2023geval} need one or more calls to a large model per summary and per dimension. Specialised reference-free metrics such as SummaC~\citep{laban2022summac} and AlignScore~\citep{zha2023alignscore} are cheap but address factual consistency only. What is missing is a metric that is cheap, needs no reference, and covers all four dimensions.

\paragraph{LNC.} We propose LNC, the Likelihood--NLI Composite, which fills this gap with two small open models. A 1.5B-parameter causal language model reads the source and the summary and provides three kinds of evidence: how plausible the summary is given the source, how much more plausible it is than minimally perturbed versions of itself (with swapped sentences, flipped negations, exchanged entities or changed numbers), and how plausible its words are at the points where it departs from the source. An NLI model checks whether every summary sentence is entailed by a nearby passage of the source. A per-dimension ridge regression, fitted on 50 SummEval articles and frozen afterwards, combines these ten signals into four scores. Every signal has prior art; what is new is the selection of signals that remain informative once surface overlap is controlled for, their per-dimension combination, and the evidence on where such a calibrated metric does and does not transfer.

\paragraph{Measuring more than overlap.} Evaluating a metric of this kind requires care, because summarisation benchmarks offer an extractive shortcut. The median SummEval summary copies 92\% of its token bigrams from the source, and the human ratings follow the amount of copying: summaries that copy more are rated as more consistent, more fluent and, to a lesser degree, more coherent and relevant. Any metric that rewards overlap with the source therefore inherits a correlation it did not earn by assessing quality. The effect is large: the bigram copy rate alone, with no model, no reference and no parameter, correlates with the human ratings better than ten of the fourteen established metrics in our pool. \citet{durmus2022spurious} and \citet{ladhak2022faithful} observed a related effect for factual consistency. We turn it into an evaluation axis for all four dimensions: every correlation is reported both raw and with the copy rate partialled out, so that a metric is credited only for what it measures beyond the shortcut, and every metric is timed on the same machine, so that cost and quality can be compared directly.

\paragraph{Results in brief.} On held-out SummEval articles LNC reaches a partial correlation of .350 at 216\,ms per summary. It is ahead of every other reference-free metric, including a local LLM judge that costs 82 times more, and it lies on the cost--quality Pareto frontier of both axes; only UniEval, which uses a reference for relevance, is more accurate. A pre-registered test with the calibration frozen shows that LNC's consistency judgement transfers to FRANK, where it outperforms UniEval and the LLM judge on the CNN/DailyMail part, matches AlignScore, and ties all three on the abstractive XSum part, whereas its relevance judgement does not transfer to RoSE.

\paragraph{Contributions.}
\begin{itemize}
  \item LNC, a reference-free metric for all four dimensions of summary quality, built from a small language model and an NLI model, more accurate than every other reference-free metric we evaluate at a fraction of the cost of an LLM judge.
  \item An evaluation protocol that credits a metric only for what it measures beyond the extractive shortcut, applied to 23 metrics and controls on SummEval with every metric timed on the same hardware.
  \item A pre-registered transfer test on FRANK and RoSE with the calibration frozen, showing that the consistency judgement generalises to new systems and annotation schemes while the relevance judgement does not.
  \item Controls and negative results that bound the claim and explain the design: part of the calibration is specific to the systems it was fitted on, a ridge over existing reference-based metrics reaches the same in-domain score, and three designs that tried to cancel the shortcut inside the language model failed for an identifiable reason.
\end{itemize}

All code, per-sample scores, the frozen calibration and the scripts that render every table in this paper are publicly available, and the whole study can be regenerated with one command.
